\documentclass[10pt,a4paper]{amsart}
\usepackage[margin=19mm]{geometry}
\usepackage{amsmath,amssymb,mathtools}
\usepackage[scr=rsfs,cal=euler]{mathalfa}
\usepackage{xfrac}
\usepackage[unicode,hidelinks,bookmarks=false]{hyperref}
\newcommand{\ie}{i.e.\ }
\newcommand{\cf}{cf.\ }
\newcommand{\resp}{respectively\ }
\newcommand{\Subsection}{\subsection}
\providecommand{\nobreakdash}{\nobreak\hbox{-}}
\setlength{\emergencystretch}{2em}
\multlinegap0pt
\usepackage{mathtools}
\newtheorem{corollary}{Corollary}
\newtheorem{proposition}{Proposition}
\newcommand{\CC}{\mathbb{C}}
\title{Towards the Relative Langlands Duality for Orthosymplectic Pairs}
\author{Dor Mezer}
\date{Source: arXiv:2606.03187v1 (CC BY 4.0)}
\begin{document}
\maketitle

\noindent\emph{Source and licence.} Towards the Relative Langlands Duality for Orthosymplectic Pairs. Version arXiv:2606.03187v1 (CC BY 4.0), distributed by arXiv under the Creative Commons Attribution 4.0 International licence, http://creativecommons.org/licenses/by/4.0/, as recorded in the arXiv licence metadata for this exact version and shown on the abstract page https://arxiv.org/abs/2606.03187v1. Full licence notice: \texttt{licenses/arxiv-2606.03187-CC-BY-4.0.txt}.\par\smallskip

\noindent\textit{Editorial scope.} Counted unit: the article's proposition proposition (source0.tex lines 1314--1316). The article's complete original proof of it is delivered with it (source0.tex lines 1317--1409, one \texttt{proof} environment of 93 lines). No mathematics is written, repaired or re-derived: the statement and the proof are the article's own lines, in the article's own notation, and no retained line is substituted or re-flowed.  Retained uncounted as the source-local context the proof needs: lines 935--935; lines 1301--1303; lines 1308--1310.  Nothing was omitted from any retained span, so the package carries no omission record.  No retained line contains a citation, so the wrapper carries no bibliography and no bibliography entry.  No retained line contains a figure environment, an image-inclusion command or drawing code, so no image is delivered and the package declares no asset dependency.  Editorial formatting only, in addition: an AMS \texttt{amsart} preamble that restates the article's own declarations and macros used by the retained lines, namely the environments \texttt{corollary}, \texttt{proposition} on separate counters, together with the standard packages those lines need.  The retained environments print in this document as Corollary 1, Proposition 1 in order of appearance, whereas the article prints the same items with the numbering of its own full document.  The complete native source of the article as distributed in the arXiv e-print for this exact version is bundled unchanged as \texttt{sources/201959/source0.tex}.
\subsection{A pseudo-slice}\label{sec: pseudo-slice}

\begin{equation}\label{eq: linear constraint}
    D_{(e,v_0^*)}\zeta'(0, fv_0^*)\ni \frac{d}{d\alpha}(\alpha.(x, v_0^*+\omega))|_{\alpha=0} = (0, \frac{-1}{n^2(n-1)^2}\langle \omega, fv_0^*\rangle fv_0^* + \pi_J(\omega)).
\end{equation}

\begin{corollary}\label{cor: pointwise isomorphism}
    The image of $\zeta'$ is equal to the translation by $v_0^*$ of a linear subspace $S\subset U'$, of dimension $n$. Moreover, for each $x\in \Sigma$, the restriction of $\zeta'$ to $\{x\}\times \Sigma_\mathfrak{gl}$ is an isomorphism onto $v_0^*+S$. In particular, $\zeta$ gives an isomorphism $\Sigma\xrightarrow{\sim} \Sigma_{2n}\times (v_0^*+S)$.
\end{corollary}

\begin{proposition}
    The above $S$ is equal to $\CC fv_0^*+N$, where $N:=\ker(f^{n-1})$ is as in Section \ref{sec: pseudo-slice}. That is, the image of $\zeta$ is equal to $\Sigma_{(e, f, h, v_0^*)}$.
\end{proposition}

\begin{proof}
    Let $$\epsilon := D_{(e,v_0^*)}\zeta'(0, fv_0^*)\in \CC fv_0^* + J.$$ By \eqref{eq: linear constraint} we have $S = v_0^* + \CC \epsilon + N$. We write $\epsilon = \alpha fv_0^* + \beta j$, with $j\in J$ satisfying $\langle j, j\rangle = -n^2(n-1)^2$.
    The claim of the proposition is equivalent to showing that $\beta=0$.
    
    We start by giving the proof in the case $n\geq 3$.
    Consider $$W:=\left(v_0^*+\mathrm{Span}(\{fv_0^*, f^2v_0^*, f^3v_0^*\})\right)\subset \Sigma_{(e, f, h, v_0^*)}.$$
    We claim that the image of $\{e\}\times W$ under $\zeta'$ is equal to $v_0^* + \mathrm{Span}(\epsilon, f^2v_0^*, f^3v_0^*)$.
    Indeed, let $w\in W$, and write $\zeta'((e, w))=b_0\epsilon + \sum_{i=1}^{n-1}b_i f^{i+1}v_0^*$. Let $i$ be the maximal integer such that $b_i\neq 0$.
    If $i > 2$, we have
    $$0 = \langle w, e^{2i} w\rangle = \langle \zeta'(e, w), e^{2i} \zeta'(e, w)\rangle = (-1)^i b_i^2 \langle e^if^{i+1}v_0^*, e^if^{i+1}v_0^*\rangle \neq 0.$$
    This shows that $\zeta'(\{e\}\times W)\subseteq v_0^* + \mathrm{Span}(\epsilon, f^2v_0^*, f^3v_0^*)$. Both are 3-dimensional affine spaces, and by Corollary \ref{cor: pointwise isomorphism}, $\{e\}\times W$ maps under $\zeta'$ isomorphically to its image. Thus, we get the desired equality.

    Let $w\in W, w = v_0^*+a_0fv_0^* + a_1f^2v_0^* + a_2 f^3v_0^*$, and $\zeta'(e, w) = v_0^* + b_0\epsilon + b_1f^2v_0^* + b_2f^3v_0^*$.
    Setting
    \[\begin{aligned}
        P_0(w) &= \frac{1}{2n^2(n-1)^2}\langle w, w\rangle = a_1-\frac{1}{2}a_0^2\\
        P_1(w) &= \frac{1}{2n^3(n-1)^3}\langle w, e^2 w\rangle = \frac{1}{2}n(n-1)a_1^2 - (n+1)(n-2)a_0a_2\\
        P_2(w) &= -\frac{1}{2n^4(n-1)^4(n+1)^2(n-2)^2}\langle w, e^4 w\rangle = \frac{1}{2}a_2^2\\
        P_0'(\zeta'(e, w)) &= \frac{1}{2n^2(n-1)^2}\langle \zeta'(e, w), \zeta'(e, w)\rangle = b_1-\frac{1}{2}(\alpha^2+\beta^2)b_0^2\\
        P_1'(\zeta'(e, w)) &= \frac{1}{2n^3(n-1)^3}\langle \zeta'(e, w), e^2 \zeta'(e, w)\rangle = \frac{1}{2}n(n-1)b_1^2 - (n+1)(n-2)\alpha b_0b_2\\
        P_2'(\zeta'(e, w)) &= -\frac{1}{2n^4(n-1)^4(n+1)^2(n-2)^2}\langle \zeta'(e, w), e^4 \zeta'(e, w)\rangle = \frac{1}{2}b_2^2
    \end{aligned}\]
    The Jacobian matrix of $(P_0, P_1, P_2)$ is $$\begin{pmatrix}
    -a_0&-(n+1)(n-2)a_2&0\\
    1&n(n-1)a_1&0\\
    0& -(n+1)(n-2)a_0 & a_2
    \end{pmatrix}$$
    and of $(P_0', P_1', P_2')$ is $$\begin{pmatrix}
    -(\alpha^2+\beta^2)b_0&-(n+1)(n-2)\alpha b_2&0\\
    1&n(n-1)b_1&0\\
    0& -(n+1)(n-2)\alpha b_0 & b_2
    \end{pmatrix}.$$
    If we restrict to the locus where $a_2$ is non-zero and the Jacobian matrix of $(P_0, P_1, P_2)$ is singular, i.e.
    $$a_2 = \frac{n(n-1)}{(n+1)(n-2)}a_0a_1,$$
    We get that the Jacobian matrix of $(P_0', P_1', P_2')$, must also be singular and $b_2\neq 0$ (by comparison of $P_2$ to $P_2'$), i.e.
    $$\alpha b_2 = \frac{n(n-1)}{(n+1)(n-2)}(\alpha^2+\beta^2)b_0b_1.$$
    We get that for any $a_0, a_1\in \CC$ which are non-zero, there are $b_0, b_1\in \CC\setminus\{0\}$ such that
    \[
    \begin{aligned}
        (P_0=P_0')&\ \ \ \ b_1-\frac{1}{2}(\alpha^2+\beta^2)b_0^2 = a_1-\frac{1}{2}a_0^2\\
        (P_1=P_1')&\ \ \ \ \frac{1}{2}b_1^2 - (\alpha^2+\beta^2)b_0^2b_1 = \frac{1}{2}a_1^2 - a_0^2a_1\\
        (P_2 = P_2')&\ \ \ \ (\alpha^2+\beta^2)^2b_0^2b_1^2 = \alpha^2 a_0^2a_1^2
    \end{aligned}
    \]

    Choosing $a_0=1, a_1 = 2$ we get that
    \[
    \begin{aligned}
        &b_1-\frac{1}{2}(\alpha^2+\beta^2)b_0^2 = \frac{3}{2}\\
        &\frac{1}{2}b_1 - (\alpha^2+\beta^2)b_0^2 = 0\\
        &(\alpha^2+\beta^2)^2b_0^2b_1^2 = 4\alpha^2
    \end{aligned}
    \]
    From the first two equation we get that $(\alpha^2+\beta^2)b_0^2=1, b_1=2$.
    Hence, the third equation reads $4(\alpha^2+\beta^2)=4\alpha^2$ and so $\beta=0$, as desired.

    Now we assume $n=2$.
    We denote $$\zeta'(e+f, v_0^* + a_0fv_0^* + a_1f^2v_0^*)= v_0^*+b_0\epsilon+b_1f^2v_0^*,$$
    where recall that $\epsilon = \alpha fv_0^* + \beta j$ and $j\in J, \langle j, j\rangle = -4$.
    We have
    \[\begin{aligned}
        P_0(a_0, a_1) &= \frac{1}{8}\langle v_0^* + a_0fv_0^* + a_1f^2v_0^*, v_0^* + a_0fv_0^* + a_1f^2v_0^*\rangle = a_1-\frac{1}{2}a_0^2\\
        P_1(a_0, a_1) &= \frac{1}{4}\langle (e+f)(v_0^* + a_0fv_0^* + a_1f^2v_0^*), (e+f) (v_0^* + a_0fv_0^* + a_1f^2v_0^*)\rangle = -(1+2a_1)^2 + 4a_0^2\\
        P_0'(b_0, b_1) &= \frac{1}{8}\langle v_0^*+b_0\epsilon+b_1f^2v_0^*, v_0^*+b_0\epsilon+b_1f^2v_0^*\rangle = b_1-\frac{1}{2}(\alpha^2+\beta^2)b_0^2\\
        P_1'(b_0, b_1) &= \frac{1}{4}\langle (e+f)(v_0^* + b_0\epsilon + b_1f^2v_0^*), (e+f) (v_0^* + b_0\epsilon + b_1f^2v_0^*)\rangle = -(1+2b_1)^2 + 4\alpha^2b_0^2\\
    \end{aligned}\]
    with $P_0 = P_0'$ and $P_1 = P_1'$.
    The Jacobian matrix of $(P_0, P_1)$ is $$\begin{pmatrix}
        -a_0 & 8a_0\\
        1 & -4(1+2a_1)
    \end{pmatrix}$$
    and its determinant is equal to $4a_0(-1+2a_1)$.
    The Jacobian matrix of $(P_0', P_1')$ is $$\begin{pmatrix}
        -b_0(\alpha^2+\beta^2)&8\alpha^2b_0\\
        1& -4(1+2b_1)
    \end{pmatrix}$$
    and its determinant is equal to $4b_0(-\alpha^2+\beta^2+(\alpha^2+\beta^2)2b_1)$.
    Setting $a_1 = \frac{1}{2}$, we get that we must have:
    \[
    \begin{aligned}
        (P_0=P_0')&\ \ \ \ b_1-\frac{1}{2}(\alpha^2+\beta^2)b_0^2 = \frac{1}{2}-\frac{1}{2}a_0^2\\
        (P_1=P_1')&\ \ \ \ -(1+2b_1)^2+4\alpha^2b_0^2 = -4+4a_0^2\\
        (\mathrm{Jacobian}=0)&\ \ \ \ 4b_0(-\alpha^2+\beta^2+(\alpha^2+\beta^2)2b_1)=0
    \end{aligned}
    \]
    From the first two equations we get
    $$(1-2b_1)^2 + 4\beta^2b_0^2 = 0$$
    We see that if $b_0=0$ then $b_1=\frac{1}{2}$, which is impossible if $a_0\neq 0$. So we must have by the third equation above
    $$2(\alpha^2+\beta^2)b_1 = \alpha^2-\beta^2.$$
    This implies that
    $$\left(\frac{\beta^2-\alpha^2}{\alpha^2+\beta^2}\right)^2 + 4\beta^2b_0^2 = 0.$$
    If $\beta\neq 0$, it would follow that the values of $b_0, b_1$ are fixed, but this is clearly impossible as $a_0$ is allowed to vary.
\end{proof}

\end{document}
